\section{State of the Art}
\label{sec:stateoftheart}

\subsection{Service-Oriented Architectures and Security}

The decomposition of monolithic applications into microservices
increases operational flexibility but multiplies the attack surface: each service endpoint
becomes a potential entry point, inter-service calls transit a shared network that may be
partially compromised, and credentials or tokens that authenticate service-to-service
communication must be protected across many processes and containers.

\subsection{Monitoring Tools}

Prometheus is an open-source, pull-based metrics platform originally developed
at SoundCloud. Its time-series data model and flexible PromQL query
language make it well-suited to microservice observability. Grafana
provides dashboarding over Prometheus and other data sources. Commercial alternatives
such as Datadog and Zabbix offer similar functionality but introduce cost and vendor
dependency. Open-source tooling is preferred in this project to maintain transparency
over the detection pipeline.

\subsection{SIEM and Intrusion Detection}

Wazuh is an open-source SIEM derived from OSSEC. It ingests
log files, applies pattern-matching and frequency rules, and can trigger active responses
(shell scripts) on rule matches. It is the platform adopted in this project for
consistency with the course practicals. Commercial platforms such as Splunk provide
richer analytics but hide the underlying rule machinery. Falco targets
Kubernetes-specific syscall monitoring and does not apply to the Incus/Docker
deployment used here.

\subsection{DDoS Attacks and Mitigation}

DDoS attacks aim to exhaust the resources of a target service, rendering it unavailable
to legitimate users. Rate-based detection (triggering when a
source IP exceeds a request threshold within a time window) is a standard first-line
defence. Its limitation is that distributed floods (many source IPs) evade single-IP
thresholds, a limitation acknowledged in the evaluation.

\subsection{Advanced Persistent Threats}

APTs are characterised by stealth, persistence, and multi-stage execution.
The reconnaissance phase (endpoint enumeration, banner grabbing) precedes exploitation.
Honeypot-based detection (placing intentionally exposed
"decoy" endpoints that log all access) is a lightweight technique that generates high-fidelity
alerts: any access to a honeypot endpoint is inherently anomalous.

\subsection{Spam and Content Filtering}

Email spam remains a persistent threat vector, particularly in alert-driven operational
environments where operators are conditioned to act on incoming alerts. SpamAssassin applies a scoring model that combines header analysis, content
rules, and Bayesian classification. Its extensible rule system allows custom scoring for
domain-specific content (e.g., pharmaceutical keywords in a city monitoring context).
Application-layer alert flooding (using legitimate credentials to trigger high volumes of
internal alerts) is a distinct attack class that content filters cannot address alone and
requires circuit-breaker or rate-limiting logic at the application level.
